\chapter{Amniocentesis}

\section*{What Is This Test or Procedure?}
Amniocentesis is a prenatal test where a small amount of amniotic fluid is removed from the sac surrounding the fetus to test for genetic abnormalities or lung maturity.

\section*{Alternative Names}
Amniotic fluid test, Genetic amniocentesis

\section*{Principle}
A thin needle is passed through the maternal abdomen into the amniotic sac under ultrasound guidance and a small volume of amniotic fluid is aspirated. The fluid contains fetal cells and biochemical products that are cultured and analyzed for chromosomes, gene mutations, proteins, and other markers of fetal health.
\section*{History}
First performed in the late 19th century to relieve polyhydramnios. Evolved in the 1970s with ultrasound guidance and cytogenetic culture.

\section*{Why This Test or Procedure Is Ordered}
Ordered for pregnant women with an increased risk of fetal genetic abnormalities (due to age or abnormal screening tests) or to check fetal lung maturity.

\section*{Is This a Primary or Secondary Test or Procedure}
Primary prenatal diagnostic test for chromosomal abnormalities.

\section*{How This Test or Procedure Is Performed}
Under ultrasound guidance, a long, thin needle is inserted through the abdomen into the amniotic sac to withdraw about 20 mL of amniotic fluid.

\section*{What Preparation Is Needed}
Signed informed consent; NPO is not required; full or empty bladder depending on gestational age.

\section*{Contraindications and Precautions}
None in high-risk pregnancies; relative includes maternal blood-borne infections (HIV, Hepatitis B/C) due to transmission risk.

\section*{Risks}
Miscarriage (0.1-0.3\%), leakage of amniotic fluid, cramping, or uterine infection.

\section*{What Happens After the Test or Procedure}
Rest for the remainder of the day; avoid heavy lifting or strenuous activity for 24-48 hours.

\section*{Understanding Your Results}
Cells from the fluid are cultured and analyzed for chromosomal karyotype, microarrays, or specific gene mutations.

\section*{Normal Results}
Normal fetal karyotype (46,XX or 46,XY); no chromosomal abnormalities detected.

\section*{Follow-up Tests and Procedures}
Genetic counseling; prenatal care planning.

\section*{References}
\begin{enumerate}
  \item MedlinePlus. Amniocentesis. U.S. National Library of Medicine; 2025.
  \item Current Medical Diagnosis and Treatment. McGraw-Hill; 2025.
\end{enumerate}

\clearpage
